\documentclass[final]{beamer}

\usepackage[orientation=portrait,size=a0,scale=1.4]{beamerposter}
\usepackage[utf8]{inputenc}
\usepackage[T1]{fontenc}
\usepackage[spanish]{babel}
\usepackage{amsmath,amssymb,amsfonts}
\usepackage{graphicx}
\usepackage{multicol}
\usepackage{lmodern}

\setbeamertemplate{navigation symbols}{}
\setbeamertemplate{footline}{}
\setbeamertemplate{blocks}[rounded][shadow=true]

\setbeamercolor{block title}{bg=blue!15,fg=black}
\setbeamercolor{block body}{bg=white,fg=black}
\setbeamerfont{block title}{size=\large,series=\bfseries}
\setbeamerfont{block body}{size=\normalsize}

\newcommand{\Prob}{\mathbb{P}}

\begin{document}

\begin{frame}[t]

% ================= PORTADA =================
\begin{center}
{\VeryHuge\bfseries Ergodicidad y convergencia en redes espacio-temporales\par}
\vspace{0.35cm}
{\Large VI Congreso Internacional Multidisciplinario de Matemática\par}
\vspace{0.35cm}
{\large Luis Andre Garro Hernandez, Aldo Abel Valeriano Macuri, Dire Daniel Espinal Pecho, Arley Ricardo Gutarra Inga\par}
{\large Universidad Nacional de Ingeniería\par}
{\large \texttt{luis.garro.contacto@gmail.com}\par}
\end{center}

\vspace{0.5cm}

\begin{columns}[t]

% ================= COLUMNA IZQUIERDA =================
\begin{column}{0.48\paperwidth}

\begin{block}{01 Introducción y objetivo}
\textbf{Problema.} Dado un grafo temporal \(G_t=(V,E_t)\), donde todo \(v_i\in V(G)\) tiene una posición \(p_i(t)\in V(X)\) y pertenece a un conjunto de color \(C_t\subseteq V(G)\), ¿cómo influyen la dinámica temporal de los enlaces en el comportamiento del sistema a largo plazo?

\vspace{0.3cm}
\textbf{Objetivo.} Mostrar que, bajo ciertas condiciones, el modelo converge a una única distribución estacionaria \(\pi^*\).
\end{block}

\begin{block}{02 Marco teórico}
\textbf{Definiciones y notación.}

\vspace{0.2cm}
\textbf{Conexión y desconexión.} Para todo instante \(t\) y par \(v_i,v_j\in V(G)\):
\[
\{v_i,v_j\}\in E(G_t)\cap \{p_i(t),p_j(t)\}\in E(X)
\;\Longrightarrow\;
\Prob\bigl(\{v_i,v_j\}\in E(G_t)\bigr)
=
\frac{(d_i^-(v_j)+1)(d_i^-(v_j)+1)}{(n+1)^2}.
\]
\[
\{v_i,v_j\}\in E(G_t)\cap \{p_i(t),p_j(t)\}\notin E(X)
\;\Longrightarrow\;
\Prob\bigl(\{v_i,v_j\}\in E(G_t)\bigr)
=
1-\frac{d_i^-(v_j)d_i^-(v_j)}{n^2}.
\]

\vspace{0.2cm}
\textbf{Coloración.} \(N_t,N_v\in V(G)\):
\[
\Prob\left[v\in \bigcup_{i=0}^{c}C_{t_i}\right]=1,
\qquad
\Prob\left[v\in C_{t_i}\cap C_{t_j}\right]=0,
\quad \forall i\ne j.
\]

\vspace{0.2cm}
\textbf{Con pesos y probabilidad de transición:}
\[
w_j(t)=|N^c_G(v)\cap C_{t_i}|+\alpha\,\mathbf{1}_{\{p_v\in C_{t_i}\}},
\]
\[
\Prob[v\in C_{t_i}]
=
k\,\mathbf{1}_{\{\nu=0\}}+(1-k)w_j(t),
\quad v\in C_{t_i}.
\]

\vspace{0.2cm}
\textbf{Desplazamiento.} En \(t^+\to(t+1)^-\), cada \(v_i\in V(G)\) elige uniformemente una posición destino
\[
x_j\in N_X(p_{t_i}(v))\cup\{p_{t_i}(v)\}
\]
y aplica la transposición espacial
\[
\sigma_{i,t}=(p_{t_i}(v_j)\;x_j)\in S_m.
\]
La nueva configuración resulta de la composición ordenada
\[
\sigma_t=\sigma_{n,t}\circ\cdots\circ\sigma_{1,t}.
\]

\vspace{0.2cm}
\textbf{Hipótesis.} El croquis \(X\) es finito, conexo \((X\ne K_m)\), \(k\in(0,1)\) y \(\alpha>0\) con \(C_i^X\ne\varnothing\) para todo \(i\in\{1,\dots,c\}\).

\vspace{0.2cm}
\textbf{Antecedentes.} Los grafos temporales discretos se modelan mediante secuencias de instantáneas [1]; los modelos coevolutivos analizan el consenso cuando red y opiniones cambian juntas [2].
\end{block}

\begin{block}{03 Método o demostración}
El comportamiento asintótico del sistema queda unívocamente determinado por la \(4\)-tupla
\[
(X,\{C_i^X\},k,\alpha),
\]
garantizando un equilibrio estadístico único independiente de las condiciones iniciales.

\vspace{0.3cm}
La demostración se divide en tres partes:
\begin{enumerate}
\item \textbf{Cadena de Markov finita:} se muestra que el conjunto de estados posibles \(\Omega\) es finito y que se cumple la propiedad de pérdida de memoria.
\item \textbf{Irreducibilidad:} se define explícitamente una secuencia finita de pasos para llevar un estado \(S_1\) a otro \(S_2\).
\item \textbf{Aperiodicidad:} como la probabilidad de que un estado \(S\) no cambie es positiva, se concluye la aperiodicidad.
\end{enumerate}
\end{block}

\end{column}

% ================= COLUMNA DERECHA =================
\begin{column}{0.48\paperwidth}

\begin{block}{04 Resultados principales}
\textbf{Teorema (Ergodicidad y equilibrio estadístico único).}
Si el croquis \(X\) es finito, conexo \((X\ne K_m)\), \(k\in(0,1)\) y \(\alpha>0\) con \(C_i^X\ne\varnothing\) para todo \(i\in\{1,\dots,c\}\), entonces \(\{D_t\}_{t\ge 0}\) es una cadena de Markov finita, irreducible y aperiódica. Por lo tanto, existe una única distribución estacionaria \(\pi^*\) tal que, para cualquier estado inicial \(\pi_0\),
\[
\lim_{t\to\infty}\pi_0P^t=\pi^*.
\]

\vspace{0.3cm}
\textbf{Interpretación.} El equilibrio alcanzado es de naturaleza estadística y no estática: el diagrama nunca deja de evolucionar, pero la frecuencia con la que visita cada configuración \(S\in\Omega\) converge a una distribución fija \(\pi^*\), completamente independiente de la red y coloración iniciales.
\end{block}

\begin{figure}
\centering
\IfFileExists{figura_visualizacion.png}{%
  \includegraphics[width=0.95\linewidth]{figura_visualizacion.png}%
}{%
  \fbox{\parbox{0.9\linewidth}{\centering
  Coloque aquí el archivo \texttt{figura\_visualizacion.png} con la visualización del póster.}}
}
\caption{Visualización con \(X=C_6\) y los estados iniciales \(G^{(1)}=K_6\), \(G^{(2)}=K_{1,5}\) y \(G^{(3)}=K_6\), con coloraciones unánimes \(C_1\), \(C_2\) y \(C_6\), respectivamente.}
\end{figure}

\begin{block}{05 Conclusiones}
Los grafos temporales discretos modelan mediante secuencias de instantáneas [1], mientras que los modelos coevolutivos analizan el consenso cuando red y opiniones cambian juntas [2].

\vspace{0.2cm}
El comportamiento asintótico del sistema queda unívocamente determinado por la \(4\)-tupla \((X,\{C_i^X\},k,\alpha)\), garantizando un equilibrio estadístico único independiente de las condiciones iniciales.

\vspace{0.2cm}
\textbf{Trabajo futuro.} Estudiar el problema inverso al fijar \((X,\{C_i^X\})\): determinar si distintos pares \((k,\alpha)\) pueden inducir una misma distribución estacionaria \(\pi^*\) y caracterizar la geometría de dichas curvas de nivel en el plano de parámetros.
\end{block}

\begin{block}{06 Referencias}
\begin{enumerate}
\item D. Wu, Y. Xu, X. Lin, W. Zhang, and Y. Zhang, \emph{Understanding Evolving Graph Structures for Large Discrete-Time Dynamic Graph Representation}, Proc. VLDB Endow., 19(5), 862--875 (2026).
\item N. Djurdjevac Conrad, N. Q. Vu, and S. Nagel, \emph{Co-evolving networks for opinion and social dynamics in agent-based models}, Chaos, vol. 34, no. 9, p. 093116, 2024. arXiv:2407.00145 (2024).
\end{enumerate}
\end{block}

\end{column}

\end{columns}

\end{frame}

\end{document}